\documentclass[9pt]{article}
\usepackage[margin=0.75in, top=0.65in, bottom=0.65in]{geometry}
\usepackage{amsmath}
\usepackage{amssymb}
\usepackage{xcolor}
\usepackage{enumitem}
\usepackage{microtype}
\usepackage{booktabs}
\usepackage{parskip}
\usepackage{titlesec}
\usepackage{multicol}

\definecolor{darkblue}{RGB}{20,55,120}
\definecolor{rulegray}{RGB}{180,180,190}

\titleformat{\section}[block]
  {\normalfont\normalsize\bfseries\color{darkblue}}
  {}{0pt}{}[\vspace{1pt}\textcolor{rulegray}{\hrule height 0.4pt}\vspace{3pt}]
\titlespacing{\section}{0pt}{7pt}{0pt}

\titleformat{\subsection}[runin]
  {\normalfont\small\bfseries}
  {}{0pt}{}[.\enspace]
\titlespacing{\subsection}{0pt}{3pt}{0pt}

\setlist[itemize]{leftmargin=1.4em, topsep=1pt, itemsep=0pt, parsep=0pt}
\setlist[enumerate]{leftmargin=1.6em, topsep=1pt, itemsep=1pt, parsep=0pt}
\setlist[description]{leftmargin=0pt, topsep=1pt, itemsep=2pt}

\pagestyle{plain}
\setlength{\parskip}{2pt}
\setlength{\parindent}{0pt}

\begin{document}


%=============================================================
\section{Average Memory Access Time (AMAT)}

\textbf{General:}\enspace $\text{AMAT} = \text{HT} + \text{MR} \times \text{MP}$\enspace
where HT = Hit Time, MR = Miss Rate, MP = Miss Penalty.

\vspace{3pt}
\begin{multicols}{2}

\subsection{Look-Aside: L1 + MM}
CPU sends requests to \textbf{both} L1 and MM simultaneously; MM is cancelled on an L1 hit. MP is only the \emph{extra} time beyond L1.
$$\text{AMAT} = \text{HR}_{L1} \times \text{HT}_{L1} + \text{MR}_{L1} \times \text{HT}_{MM}$$

\subsection{Look-Through: L1 + MM}
CPU checks L1 first; accesses MM only on a miss, paying the full MM hit time as the penalty.
$$\text{AMAT} = \text{HT}_{L1} + \text{MR}_{L1} \times \text{HT}_{MM}$$

\columnbreak

\subsection{Look-Aside: L1 + L2 + MM}
L1 and L2 checked simultaneously; MM accessed only if both miss.
$$\text{AMAT} = \text{HR}_{L1} \cdot \text{HT}_{L1} + \text{MR}_{L1}(\text{HR}_{L2} \cdot \text{HT}_{L2} + \text{HT}_{MM} \cdot \text{MR}_{L2})$$

\subsection{Look-Through: L1 + L2 + MM}
Sequential: check L1, then L2 on miss, then MM on L2 miss.
$$\text{AMAT} = \text{HT}_{L1} + \text{MR}_{L1}\bigl[\text{HT}_{L2} + \text{MR}_{L2} \cdot \text{HT}_{MM}\bigr]$$

\end{multicols}

%=============================================================
\section{Cache Designs}

\begin{multicols}{2}

\begin{description}
  \item[\textbf{Fully Associative.}] A block can go in \emph{any} cache line---maximum flexibility, but all tags must be checked in parallel (costly hardware).
  \item[\textbf{Direct Mapped.}] Each block maps to \emph{exactly one} line (address mod \#lines)---fast and simple, but prone to conflict misses.
  \item[\textbf{Set Associative} ($k$-way)\textbf{.}] Cache split into sets; a block maps to one set and can occupy any of its $k$ lines---a practical balance.
\end{description}

\columnbreak

\textbf{Address decomposition} ($M$ = addr bits, $S$ = \#sets, $B$ = block bytes):

\vspace{3pt}
\begin{tabular}{@{}lll@{}}
  \toprule
  \textbf{Field} & \textbf{Bits} & \textbf{Note} \\
  \midrule
  Offset & $o = \log_2 B$   & byte within block \\
  Index  & $i = \log_2 S$   & selects set; 0 if fully assoc. \\
  Tag    & $t = M - i - o$  & identifies the block \\
  \bottomrule
\end{tabular}

\vspace{3pt}
\textbf{Total cache bits:}\enspace $S \times k \times (8B + t + 1)$\enspace{\small(+1 valid bit/line)}
\textbf{Order:}\enspace [$t$] [$i$] [$o$] in the memory address binary

\end{multicols}

%=============================================================
\section{Branch Prediction}

\begin{multicols}{2}

\subsection{Static vs.\ Dynamic}
\textbf{Static:} decided at compile time (e.g., always not-taken, backward branches taken); does not adapt to runtime behavior.
\textbf{Dynamic:} hardware tracks branch outcomes at runtime via a Branch Prediction Buffer and updates predictions accordingly.

\subsection{2-Level Prediction Scheme}
A \textbf{Branch History Register (BHR)} shift-records the last $k$ branch outcomes (T/NT). Its value indexes a \textbf{Pattern History Table (PHT)} of 2-bit counters, capturing inter-branch correlations for higher accuracy.

\columnbreak

\subsection{2-Bit Saturating Counter}
Requires \textbf{two} consecutive mispredictions to change the prediction, tolerating single anomalies.

\vspace{3pt}
\begin{tabular}{@{}cccc@{}}
  \toprule
  \textbf{State} & \textbf{Prediction} & \textbf{Taken} & \textbf{Not-Taken} \\
  \midrule
  \texttt{11} & Strongly Taken     & \texttt{11} & \texttt{10} \\
  \texttt{10} & Weakly Taken       & \texttt{11} & \texttt{01} \\
  \texttt{01} & Weakly Not-Taken   & \texttt{10} & \texttt{00} \\
  \texttt{00} & Strongly Not-Taken & \texttt{01} & \texttt{00} \\
  \bottomrule
\end{tabular}

\end{multicols}

%=============================================================
\section{Single-Cycle vs.\ Pipelined Execution}

\begin{multicols}{2}

\textbf{Single-cycle:} clock period = worst-case instruction path; every instruction takes exactly one (long) cycle:
$$T_{\text{clk}} = \max(\text{all instruction delays})$$

\textbf{Pipelined:} clock period = slowest stage; one instruction completes \emph{every} cycle at steady state:
$$T_{\text{clk}} = \max(\text{stage delays})\qquad \text{speedup} \approx N$$

\columnbreak

\subsection{Cycle Count ($N$ stages, $I$ instructions, $s$ stalls)}
$$\boxed{\text{Total cycles} = N + (I-1) + s}$$
The first instruction drains in $N$ cycles; each additional instruction adds 1; each stall bubble adds 1. With no hazards ($s=0$): cycles $= N + I - 1$.

\end{multicols}

%=============================================================
\section{5-Stage Pipeline}

\begin{enumerate}
  \item \textbf{IF --- Instruction Fetch.} Fetch the instruction from memory using the PC; increment the PC.
  \item \textbf{ID --- Instruction Decode / Register Read.} Decode the opcode, generate control signals, read source register values, and evaluate branch conditions.
  \item \textbf{EX --- Execute.} ALU performs the operation or computes a memory address.
  \item \textbf{MEM --- Memory Access.} Loads read from data memory; stores write to it; other instructions pass through.
  \item \textbf{WB --- Write Back.} The ALU result or loaded value is written to the destination register.
\end{enumerate}

%=============================================================
\section{Pipeline Hazards}

\begin{multicols}{2}

\subsection{Structural Hazards}
\textit{Cause:} Two instructions need the same hardware resource simultaneously (e.g., one memory for both IF and MEM).
\textit{Solutions:} \textbf{stall} until free; or \textbf{duplicate hardware} (separate I-cache and D-cache, extra ALUs).

\vspace{4pt}
\subsection{Data Hazards}
\textit{Cause:} An instruction needs a value not yet written by a prior instruction in the pipeline (\textbf{RAW} is most common).
\textit{Solutions:}
\begin{itemize}
  \item \textbf{Forwarding/Bypassing:} Route the result from EX or MEM/WB back to EX---eliminates most RAW stalls.
  \item \textbf{Load-Use Stall:} Load result unavailable until after MEM; insert one bubble (forwarding insufficient).
  \item \textbf{Code Reordering:} Compiler fills hazard slots with independent instructions.
\end{itemize}

\columnbreak

\subsection{Control Hazards}
\textit{Cause:} Instructions are fetched after a branch before its outcome is known---potentially the wrong path.
\textit{Solutions:}
\begin{itemize}
  \item \textbf{Stall/Flush:} Wait for branch resolution; flush wrong instructions. Simple but costly.
  \item \textbf{Branch Prediction:} Speculatively fetch the predicted path; flush only on misprediction.
  \item \textbf{Delayed Branching:} Always execute the instruction(s) in the \emph{delay slot} after the branch; compiler fills slots with useful independent work.
\end{itemize}

\end{multicols}

%=============================================================
\section{Forwarding Rules}

\textit{Note: Branch instructions consume their source values in the \textbf{ID} stage. R-type instructions and load/store address calculations consume values in the \textbf{EX} stage.}

\vspace{4pt}
\begin{tabular}{@{}llll@{}}
  \toprule
  \textbf{Producer} & \textbf{Consumer} & \textbf{Forward Path} & \textbf{Stalls} \\
  \midrule
  R-type & R-type or LW/SW (addr) & EX $\to$ EX & 0 \\
  R-type & Branch                 & EX $\to$ ID & 1 \\
  \addlinespace
  LW     & R-type                 & MEM $\to$ EX & 1 \\
  LW     & Branch                 & MEM $\to$ ID & 2 \\
  LW     & LW/SW (addr)           & MEM $\to$ EX & 1 \\
  \bottomrule
\end{tabular}

\vspace{6pt}
\textbf{Branch misprediction:} When a branch exits the ID stage and the prediction is wrong, the instruction fetched behind it is turned into a stall (bubble) and the correct target instruction is loaded in its place.

\end{document}
